\documentclass[11pt]{article}
\usepackage[margin=2.5cm]{geometry}
\usepackage{amsmath,amssymb}
\usepackage{graphicx}
\usepackage{hyperref}

\title{Deriving Feeding Rate from a DEB Model}
\author{}
\date{}

\begin{document}

\maketitle

\section{Goal}

This document explains how to derive a \emph{feeding rate} $f_{\text{feed}}$ (number of feeding motions per day) step-by-step from a Dynamic Energy Budget (DEB) model, given:
\begin{itemize}
  \item A species-specific parameter set ($\kappa_X$, $\{\dot{p}_{Am}\}$)
  \item Environmental food availability ($X$, $\mu_X$)
  \item Body size ($L$, $V$)
\end{itemize}

It is then the task of the behavioural control system to allocate these feeding motions within the day, subject to constraints such as bouts of repetitive feeding motions and pauses. In practice, feeding can be described as an oscillatory process with a (possibly time-varying) instantaneous frequency that integrates over 24 h to the total number $f_{\text{feed}}$.

Core parameters, state variables, and environmental quantities are listed in Appendix~\ref{app:params}.

\section{Assumptions}

We adopt two assumptions:
\begin{enumerate}
  \item We focus on \emph{rod-like isomorph larvae}, where uptake surfaces scale with $L^2$ under constant shape coefficients (Appendix~\ref{app:rod}).
  \item Each feeding motion ingests a \emph{constant fraction} of body volume/area (to be decided which assumption is more realistic): $V_{\text{bite}} = \kappa_{\text{bite}} \cdot V$ or $V_{\text{bite}} = \kappa_{\text{bite}} \cdot L^2$.
\end{enumerate}

\section{Feeding frequency computation}

\subsection*{Consumed food volume flux}

Link between food and assimilation energy fluxes:
\begin{equation}
  \dot{p}_A = \kappa_X\,\dot{p}_X\tag{1}
\end{equation}

Food volume flux follows from food energy flux:
\begin{equation}
  \dot{V}_X = \frac{\dot{p}_X}{\mu_X\,X}\tag{2}
\end{equation}

Combining these:
\begin{equation}
  \dot{V}_X = \frac{\dot{p}_A}{\kappa_X\,\mu_X\,X}\tag{3}
\end{equation}

Unit check for Eq.~(3):
\begin{equation*}
  [\dot{V}_X] = \frac{\text{energy}/\text{time}}{(\text{energy}/\text{C-mol})(\text{C-mol}/\text{volume})} = \text{volume}/\text{time}
\end{equation*}

\subsection*{Functional response and $K$}

\begin{equation}
  f(X) = \frac{X}{K + X}, \quad 0 \le f \le 1\tag{4}
\end{equation}

For an isomorph, assimilation scales with $L^2$:
\begin{equation}
  \dot{p}_A = f\,\{\dot{p}_{Am}\}\,L^2\tag{5}
\end{equation}

Assimilation capacity relates to maximum ingestion as
\begin{equation}
  \{\dot{p}_{Am}\} = \kappa_X\,\mu_X\,\{J_{Xm}\}\tag{6}
\end{equation}

This yields
\begin{equation}
  K = \frac{\{J_{Xm}\}}{\{\dot{F}_m\}}\tag{7}
\end{equation}

\subsection*{Substitution of $\dot p_A$ into Eq.~(3)}

Start from Eq.~(3) and substitute $\dot p_A$ using Eqs.~(4)--(5):
\begin{equation}\tag{8}
  \dot V_X = \frac{\dot p_A}{\kappa_X\mu_X X} = \frac{f\,\{\dot p_{Am}\}\,L^2}{\kappa_X\mu_X X}
\end{equation}

Substitute $f(X)$ from Eq.~(4) (the factor $X$ cancels):
\begin{equation}\tag{9}
  \dot V_X = \frac{\{\dot p_{Am}\}\,L^2}{\kappa_X\mu_X X}\,\frac{X}{K+X} = \frac{\{\dot p_{Am}\}\,L^2}{\kappa_X\mu_X\,(K+X)}
\end{equation}

If desired, eliminate $K$ using Eqs.~(6)--(7). From Eq.~(6),
\begin{equation}
  \{J_{Xm}\} = \frac{\{\dot p_{Am}\}}{\kappa_X\mu_X}\tag{10}
\end{equation}
so Eq.~(7) gives
\begin{equation}
  K = \frac{\{J_{Xm}\}}{\{\dot F_m\}}
    = \frac{\{\dot p_{Am}\}}{\kappa_X\mu_X\,\{\dot F_m\}}\tag{11}
\end{equation}

Substituting Eq.~(11) into Eq.~(9) yields an explicit expression for the consumed food volume flux as a function of $\{\dot p_{Am}\}$ and $X$ (with the remaining DEB constants shown):
\begin{equation}\tag{12}
  \dot V_X = \frac{\{\dot p_{Am}\}\,L^2}{\kappa_X\mu_X\left(X + \frac{\{\dot p_{Am}\}}{\kappa_X\mu_X\,\{\dot F_m\}}\right)} = \frac{\{\dot p_{Am}\}\,L^2}{\kappa_X\mu_X X + \frac{\{\dot p_{Am}\}}{\{\dot F_m\}}}
\end{equation}

Unit check for Eq.~(12):
\begin{equation*}
  [\dot V_X] = \frac{[\{\dot p_{Am}\}]\,[L^2]}{[\kappa_X\mu_X X] + \left[\frac{\{\dot p_{Am}\}}{\{\dot F_m\}}\right]} = \frac{(\text{energy}\,\text{area}^{-1}\,\text{time}^{-1})(\text{area})}{\underbrace{(\text{energy}\,\text{C-mol}^{-1})(\text{C-mol}\,\text{volume}^{-1})}_{\kappa_X\mu_X X\,\sim\,\text{energy}/\text{volume}} + \underbrace{\frac{(\text{energy}\,\text{area}^{-1}\,\text{time}^{-1})}{(\text{volume}\,\text{area}^{-1}\,\text{time}^{-1})}}_{\{\dot p_{Am}\}/\{\dot F_m\}\,\sim\,\text{energy}/\text{volume}}} = \text{volume}\,\text{time}^{-1}
\end{equation*}

\subsection*{Number of feeding motions required}

We define the feeding motion rate (feeding motions per unit time) as consumed food-volume flux divided by the volume per motion:
\begin{equation}
  f_{\text{feed}} = \dot N_{bite} = \frac{\dot V_X}{V_{\text{bite}}}\tag{13}
\end{equation}

Two alternative assumptions for $V_{\text{bite}}$ are commonly considered in applications (to be decided empirically).

\paragraph{Alternative A (volume-fraction bite).} Each feeding motion ingests a constant fraction of body volume $V$:
\begin{equation}
  V_{\text{bite}} = \kappa_{\text{bite},V}\,V = \kappa_{\text{bite},V}\,L^3\tag{14}
\end{equation}

Combining Eqs.~(13)--(14):
\begin{equation}
  f_{\text{feed}} = \frac{\dot V_X}{\kappa_{\text{bite},V}\,V} = \frac{\dot V_X}{\kappa_{\text{bite},V}\,L^3}\tag{15}
\end{equation}

Substituting $\dot V_X$ from Eq.~(12) gives an explicit form in terms of $\{\dot p_{Am}\}$ and $X$:
\begin{equation}
  f_{\text{feed}}
  = \frac{1}{\kappa_{\text{bite},V}\,L^3}\,\frac{\{\dot p_{Am}\}\,L^2}{\kappa_X\mu_X X + \frac{\{\dot p_{Am}\}}{\{\dot F_m\}}}
  = \frac{\{\dot p_{Am}\}}{\kappa_{\text{bite},V}\,L\left(\kappa_X\mu_X X + \frac{\{\dot p_{Am}\}}{\{\dot F_m\}}\right)}\tag{16}
\end{equation}

Unit check for Eq.~(16):
\begin{equation*}
  [f_{\text{feed}}] = \frac{[\dot V_X]}{[V_{\text{bite}}]} = \frac{\text{volume}\,\text{time}^{-1}}{[\kappa_{\text{bite},V}]\,[L^3]} = \frac{\text{volume}\,\text{time}^{-1}}{(1)\,(\text{volume})} = \text{time}^{-1}
\end{equation*}

\paragraph{Alternative B (area-proportional bite).} Each feeding motion ingests a volume proportional to $L^2$:
\begin{equation}
  V_{\text{bite}} = \kappa_{\text{bite},A}\,L^2\tag{17}
\end{equation}

Here $\kappa_{\text{bite},A}$ has units of length so that $V_{\text{bite}}$ has units of volume.

Combining Eqs.~(13) and (17):
\begin{equation}
  f_{\text{feed}} = \frac{\dot V_X}{\kappa_{\text{bite},A}\,L^2}\tag{18}
\end{equation}

Substituting $\dot V_X$ from Eq.~(12) yields a size-cancelled form:
\begin{equation}
  f_{\text{feed}}
  = \frac{1}{\kappa_{\text{bite},A}}\,\frac{\{\dot p_{Am}\}}{\kappa_X\mu_X X + \frac{\{\dot p_{Am}\}}{\{\dot F_m\}}}\tag{19}
\end{equation}

Unit check for Eq.~(19):
\begin{equation*}
  [f_{\text{feed}}] = \frac{[\dot V_X]}{[V_{\text{bite}}]} = \frac{\text{volume}\,\text{time}^{-1}}{[\kappa_{\text{bite},A}]\,[L^2]} = \frac{\text{volume}\,\text{time}^{-1}}{(\text{length})\,(\text{area})} = \frac{\text{volume}\,\text{time}^{-1}}{\text{volume}} = \text{time}^{-1}
\end{equation*}

\appendix

% Slightly increase row spacing in tables
\renewcommand{\arraystretch}{1.2}

\section{Parameters and variables}
\label{app:params}

\subsection*{Species-specific DEB parameters}

\begin{center}
\begin{tabular}{l|l|l|l}
  \hline
  Notation & Description & Units & Equation \\
  \hline\hline
  $\{\dot{p}_{Am}\}$ & max assimilation rate per area & energy\,area$^{-1}$\,time$^{-1}$ & $\{\dot{p}_{Am}\} = \kappa_X\,\mu_X\,\{\dot{J}_{Xm}\}$ \\
  $\{\dot{F}_m\}$ & max searching/clearance rate per area & volume\,area$^{-1}$\,time$^{-1}$ & -- \\
  $\{\dot{J}_{Xm}\}$ & max ingestion rate per area & C-mol\,area$^{-1}$\,time$^{-1}$ & -- \\
  $\kappa_X$ & food-to-reserve conversion efficiency & -- & $\dot{p}_A = \kappa_X\,\dot{p}_X$ \\
  $\kappa_{\text{bite},V}$ & bite volume fraction (Alt.~A) & -- & $V_{\text{bite}} = \kappa_{\text{bite},V}\,V$ \\
  $\kappa_{\text{bite},A}$ & bite area-proportional fraction (Alt.~B) & length & $V_{\text{bite}} = \kappa_{\text{bite},A}\,L^2$ \\
  $\delta_M$ & DEB shape coefficient & -- & $L = \delta_M L_w$ \\
  \hline
\end{tabular}
\end{center}

\subsection*{State variables}

\begin{center}
\begin{tabular}{l|l|l|l}
  \hline
  Notation & Description & Units & Equation \\
  \hline\hline
  $L$ & structural length & length & -- \\
  $V$ & structural volume & volume & $V = L^3$ \\
  \hline
\end{tabular}
\end{center}

\subsection*{Environmental and derived quantities}

\begin{center}
\begin{tabular}{l|l|l|l}
  \hline
  Notation & Description & Units & Equation \\
  \hline\hline
  $X$ & food concentration & C-mol\,volume$^{-1}$ & -- \\
  $\mu_X$ & energy per C-mol & energy\,C-mol$^{-1}$ & -- \\
  $K$ & half-saturation constant & C-mol\,volume$^{-1}$ & $K = \{J_{Xm}\}/\{\dot{F}_m\}$ \\
  $f$ & functional response & -- & $f = X/(K+X)$ \\
  $\dot{p}_A$ & assimilation flux & energy\,time$^{-1}$ & $\dot{p}_A = f\,\{\dot{p}_{Am}\}\,L^2$ \\
  $\dot{p}_X$ & food energy flux & energy\,time$^{-1}$ & $\dot{p}_A = \kappa_X\,\dot{p}_X$ \\
  $\dot{V}_X$ & food volume flux & volume\,time$^{-1}$ & $\dot{V}_X = \dot{p}_A/(\kappa_X\,\mu_X X)$ \\
  $V_{\text{bite}}$ & volume per feeding motion & volume & $V_{\text{bite}} = \kappa_{\text{bite},V}\,V$ or $\kappa_{\text{bite},A}\,L^2$ \\
  $f_{\text{feed}}$ & feeding motions per time & time$^{-1}$ & $f_{\text{feed}} = \dot{V}_X/V_{\text{bite}}$ \\
  \hline
\end{tabular}
\end{center}

\section{Larval body as a rod-like cylindrical isomorph}
\label{app:rod}

We treat the larval body as an \emph{isomorph} in the DEB sense: shape is constant during growth, so structural \emph{volume} scales with the cube of structural \emph{length}, and structural \emph{surface area} scales with the square of structural length.

For convenience (and because many larvae are approximately elongated), we represent the body as a \emph{right circular cylinder} with a constant aspect ratio.

Notation note: we use $L_w$ for observed (physical) length to avoid confusion with the DEB \emph{scaled} length, which is commonly written as $l$.

\subsection*{Observed (physical) dimensions}

Let $L_w$ denote the observed body length (cm) and $R$ the body radius (cm). We approximate structural volume $V$ by the geometric cylinder volume
\begin{equation}
  V = \pi R^2 L_w
\end{equation}

The total cylinder surface area (including end caps) is
\begin{equation}
  A_{\mathrm{cyl}} = 2\pi R L_w + 2\pi R^2
\end{equation}

and for elongated larvae ($L_w \gg R$) the lateral area dominates:
\begin{equation}
  A_{\mathrm{cyl}} \approx 2\pi R L_w
\end{equation}

\subsection*{Constant-shape (isomorph) constraint}

We impose a constant aspect ratio
\begin{equation}
  \alpha_{\mathrm{AR}} \equiv \frac{L_w}{R} = \mathrm{const}
\end{equation}
so $R = L_w/\alpha_{\mathrm{AR}}$ and therefore
\begin{equation}
  V = \frac{\pi}{\alpha_{\mathrm{AR}}^2}\,L_w^3
\end{equation}

\subsection*{Structural length and the DEB shape coefficient}

In DEB, structural volume and structural length are related by
\begin{equation}
  V = L^3
\end{equation}

For constant shape, observed length $L_w$ is proportional to structural length $L$ via the shape coefficient $\delta_M$ (dimensionless):
\begin{equation}
  L = \delta_M\,L_w
\end{equation}

For the cylindrical rod with constant aspect ratio $\alpha_{\mathrm{AR}}$, combining $V=L^3$ with $V=(\pi/\alpha_{\mathrm{AR}}^2)L_w^3$ gives
\begin{equation}
  \delta_M = \left(\frac{\pi}{\alpha_{\mathrm{AR}}^2}\right)^{1/3}
\end{equation}

Equivalently,
\begin{equation}
  L_w = \frac{L}{\delta_M} = \left(\frac{\alpha_{\mathrm{AR}}^2}{\pi}\right)^{1/3} L
\end{equation}

\subsection*{Implications for DEB flux scaling}

For an isomorph, DEB writes size dependence using structural length $L$ with $V=L^3$ and (implicitly) structural surface area proportional to $L^2$. Below are the basic relations used in analyses.

\paragraph{Surface-area controlled fluxes and capacities ($\propto L^2$).}

Assimilation power (energy/time) is written with the surface-area-specific maximum $\{p_{Am}\}$:
\begin{equation}
  \dot p_A = f\,\dot p_{Am} = f\,\{p_{Am}\}\,L^2
\end{equation}

Maximum assimilation power at size $L$ is
\begin{equation}
  \dot p_{Am} = \{p_{Am}\}\,L^2
\end{equation}

For feeding derived from searching/filtering (Chapter 2.1.4), the maximum searching rate scales as
\begin{equation}
  \dot F_m = \{\dot F_m\}\,L^2
\end{equation}

and the maximum ingestion flux of food (amount/time) is
\begin{equation}
  \dot J_{X m} = \{\dot J_{X m}\}\,L^2
\end{equation}

With food density $X$ and half-saturation coefficient $K$, the (hyperbolic) functional response is
\begin{equation}
  f = \frac{X}{K+X}
\end{equation}

and ingestion follows
\begin{equation}
  \dot J_X = f\,\dot J_{X m} = f\,\{\dot J_{X m}\}\,L^2
\end{equation}

The DEB textbook link between $K$, searching, and maximum ingestion is
\begin{equation}
  K = \frac{\dot J_{X m}}{\dot F_m} = \frac{\{\dot J_{X m}\}}{\{\dot F_m\}}
\end{equation}

\paragraph{Volume-controlled costs ($\propto L^3$).}

Somatic maintenance power is written using the volume-specific maintenance rate $[p_M]$:
\begin{equation}
  \dot p_M = [p_M] L^3
\end{equation}

\section{Practical note -- measuring $\kappa_{\text{bite}}$}

In future experiments, the bite parameters can be estimated by combining measurements of body size, food volume flux $\dot V_X$, and feeding frequency $f_{\text{feed}}$.

For Alternative A (Eq.~14):
\begin{equation}
  \kappa_{\text{bite},V} = \frac{\dot V_X}{f_{\text{feed}}\,V}
\end{equation}

For Alternative B (Eq.~17):
\begin{equation}
  \kappa_{\text{bite},A} = \frac{\dot V_X}{f_{\text{feed}}\,L^2}
\end{equation}

At present we treat $\kappa_{\text{bite},V}$ or $\kappa_{\text{bite},A}$ as a species- and instar-specific quantity to be quantified in dedicated feeding trials.

\section*{References}

\begin{itemize}
  \item Kooijman, S.A.L.M. (2010). \emph{Dynamic Energy Budget theory for metabolic organisation} (3rd ed.). Cambridge University Press.
  \item Standard DEB notation: \url{https://www.bio.vu.nl/thb/deb/deblab/add_my_pet/}
\end{itemize}

\end{document}
